% Zumdahl Chapter 12 free-response set -- 2 long (10) + 3 short (4) = 32 pts.
\documentclass{apchem-exam}

\apcunitword{Chapter}
\apcunit{12}
\apcunittitle{Chemical Kinetics}
\apcdoctitle{Free-Response Set}
\apcsubtitle{Zumdahl \S12.1--\S12.7 \textbullet\ 5 questions \textbullet\ 32 points \textbullet\ 50 minutes}

\begin{document}
\apctitleblock

\begin{apcnote}[Scope --- one CED exclusion applies]
Covers \textbf{CED 5.1} (\S12.1), \textbf{5.2} (\S12.2--12.3),
\textbf{5.3} (\S12.4), \textbf{5.4} and \textbf{5.7--5.9} (\S12.5),
\textbf{5.5}, \textbf{5.6} and \textbf{5.10} (\S12.6), and \textbf{5.11}
(\S12.7).

The CED exclusion on topic 5.6 removes \textbf{calculations involving the
Arrhenius equation}. You will not be asked to compute an activation energy
from two rate constants, or to use $\ln k$ against $1/T$ to extract
$E_a$. Reading an energy profile and reasoning qualitatively about how
temperature and catalysts change a rate both remain fully examinable, and
both appear below.
\end{apcnote}

\examsection{II}{Free Response}{5 questions \textbullet\ 32 points \textbullet\ 50 minutes}{\calculatorok}

% =====================================================================
\frq{long}{10}{Zumdahl \S12.2--\S12.3 \textbullet\ CED 5.2 \textbullet\ determining a rate law from initial rates}

For the reaction \ce{A + B -> C}, the following initial rates were
measured at constant temperature.

\begin{center}
\begin{tabular}{cccc}
\hline
\textbf{Experiment} & $[\ce{A}]_0$ (M) & $[\ce{B}]_0$ (M) &
\textbf{Initial rate} (M/s) \\
\hline
1 & 0.10 & 0.10 & $2.0\times10^{-3}$ \\
2 & 0.20 & 0.10 & $8.0\times10^{-3}$ \\
3 & 0.10 & 0.20 & $4.0\times10^{-3}$ \\
\hline
\end{tabular}
\end{center}

\begin{parts}
  \item Determine the order with respect to \ce{A}. \workspace[2.4cm]
        \keyonly{\begin{apcanswer}Compare experiments 1 and 2, where
        $[\ce{B}]$ is held constant. Doubling $[\ce{A}]$ multiplies the
        rate by $8.0/2.0 = 4$.

        Since $2^n = 4$, $n = \mathbf{2}$: second order in
        \ce{A}.\end{apcanswer}}
  \item Determine the order with respect to \ce{B}.
        \answerlines[2]{Compare experiments 1 and 3, with $[\ce{A}]$
        constant: doubling $[\ce{B}]$ doubles the rate, so
        $2^n = 2$ and $n = \mathbf{1}$ --- first order in \ce{B}.}
  \item Write the rate law and state the overall order.
        \answerlines[2]{$\text{rate} = k[\ce{A}]^2[\ce{B}]$; overall order
        $2 + 1 = \mathbf{3}$.}
  \item Calculate the value of $k$, with units. \workspace[2.8cm]
        \keyonly{\begin{apcanswer}Using experiment 1:

        $k = \dfrac{2.0\times10^{-3}}{(0.10)^2(0.10)} =
        \dfrac{2.0\times10^{-3}}{1.0\times10^{-3}} = \mathbf{2.0}$

        Units: rate is \si{\Molar\per\second} and $[\ce{A}]^2[\ce{B}]$ is
        \si{\Molar\cubed}, so $k$ is in
        \textbf{M\textsuperscript{-2}\,s\textsuperscript{-1}}.

        Check against experiment 2: $2.0(0.20)^2(0.10) =
        8.0\times10^{-3}$ \checkmark\end{apcanswer}}
  \item Explain why the orders cannot be read off the balanced equation.
        \answerlines[2]{Reaction orders are determined experimentally.
        They reflect the mechanism --- in particular the slowest step ---
        not the overall stoichiometry, and the two coincide only for a
        genuinely elementary reaction.}
\end{parts}

\begin{rubric}
  \rpoint{2}{(a) 1 pt identifying that \ce{B} is held constant; 1 pt
    order 2}
  \rpoint{2}{(b) 1 pt the comparison; 1 pt order 1}
  \rpoint{2}{(c) 1 pt the rate law; 1 pt overall order 3}
  \rpoint{2}{(d) 1 pt $k = \num{2.0}$; 1 pt correct units. A bare number
    with no units earns 1}
  \rpoint{2}{(e) 1 pt orders are experimental; 1 pt they follow the
    mechanism, not the coefficients}
\end{rubric}

% =====================================================================
\frq{long}{10}{Zumdahl \S12.5--\S12.7 \textbullet\ CED 5.4, 5.7--5.11 \textbullet\ mechanisms and catalysis}

A proposed two-step mechanism for \ce{2NO2 + F2 -> 2NO2F} is
\begin{align*}
\text{Step 1:}\quad & \ce{NO2 + F2 -> NO2F + F} \qquad \text{(slow)}\\
\text{Step 2:}\quad & \ce{NO2 + F -> NO2F} \qquad\qquad \text{(fast)}
\end{align*}

\begin{parts}
  \item Show that the mechanism sums to the overall equation.
        \answerlines[2]{Adding: \ce{2NO2 + F2 + F -> 2NO2F + F}. The
        fluorine atom appears on both sides and cancels, leaving
        \ce{2NO2 + F2 -> 2NO2F}. \checkmark}
  \item Identify any intermediate, and state how you recognized it.
        \answerlines[2]{\ce{F}, the fluorine atom. It is \emph{produced}
        in step 1 and \emph{consumed} in step 2, so it appears in the
        mechanism but not in the overall equation.}
  \item Write the rate law predicted by this mechanism, and justify it.
        \workspace[2.8cm]
        \keyonly{\begin{apcanswer}The overall rate can be no faster than
        its slowest step, so the rate law comes from step 1. Step 1 is
        \textbf{elementary}, so its orders \emph{may} be read from its
        own coefficients:

        $\text{rate} = k[\ce{NO2}][\ce{F2}]$

        Note this is first order in \ce{NO2}, even though the overall
        equation shows a coefficient of 2 --- which is exactly why the
        overall equation cannot be used to predict a rate
        law.\end{apcanswer}}
  \item Sketch in words how the energy profile for this two-step
        mechanism differs from that of a single-step reaction.
        \answerlines[3]{It has \textbf{two} maxima, one for each step,
        separated by a shallow minimum where the intermediate sits. The
        first maximum is the higher of the two, because step 1 is the
        slow step and therefore has the larger activation energy.}
  \item Explain how a catalyst increases the rate, and state what it does
        \emph{not} change. \workspace[2.8cm]
        \keyonly{\begin{apcanswer}A catalyst provides an \textbf{alternative
        pathway with a lower activation energy}. At a given temperature a
        larger fraction of collisions then has enough energy to react, so
        the rate rises.

        It does \textbf{not} change $\Delta H$ for the reaction, nor the
        position of the equilibrium: the reactants and products sit at the
        same energies as before. A catalyst is also regenerated, so it
        does not appear in the overall equation. It speeds the approach to
        equilibrium without altering where equilibrium
        lies.\end{apcanswer}}
\end{parts}

\begin{rubric}
  \rpoint{1}{(a) the steps added with \ce{F} cancelling}
  \rpoint{2}{(b) 1 pt \ce{F}; 1 pt produced then consumed}
  \rpoint{3}{(c) 1 pt rate law from the slow step; 1 pt
    $k[\ce{NO2}][\ce{F2}]$; 1 pt noting coefficients may be used only
    because the step is elementary}
  \rpoint{2}{(d) 1 pt two maxima with an intermediate between; 1 pt the
    first barrier is higher}
  \rpoint{2}{(e) 1 pt lower activation energy via an alternative pathway;
    1 pt $\Delta H$ and the equilibrium position are unchanged. Saying a
    catalyst ``lowers the energy of the reactants'' earns 0}
\end{rubric}

% =====================================================================
\frq{short}{4}{Zumdahl \S12.4 \textbullet\ CED 5.3 \textbullet\ first-order half-life}

A first-order decomposition has $k = \SI{0.0250}{\per\second}$.

\begin{parts}
  \item Calculate the half-life. \answerlines[2]{$t_{1/2} = \dfrac{0.693}
        {k} = \dfrac{0.693}{0.0250} = \mathbf{27.7}$~s}
  \item Starting from \SI{1.00}{\Molar}, state the concentration remaining
        after three half-lives, and explain why the half-life does not
        change as the reaction proceeds.
        \answerlines[3]{After three half-lives,
        $1.00 \to 0.500 \to 0.250 \to \mathbf{0.125}$~M.

        For a first-order reaction the half-life depends only on $k$, not
        on concentration: as the concentration falls the rate falls in
        exact proportion, so each successive halving takes the same time.}
\end{parts}

\begin{rubric}
  \rpoint{2}{(a) 1 pt the correct relationship; 1 pt \SI{27.7}{\second}}
  \rpoint{2}{(b) 1 pt \SI{0.125}{\Molar}; 1 pt half-life independent of
    concentration for first order}
\end{rubric}

% =====================================================================
\frq{short}{4}{Zumdahl \S12.6 \textbullet\ CED 5.5 \textbullet\ the collision model}

\begin{parts}
  \item State the two requirements for a collision to result in reaction.
        \answerlines[2]{Sufficient energy --- at least the activation
        energy --- and correct \textbf{orientation} of the colliding
        molecules.}
  \item A modest rise in temperature can double a reaction rate, even
        though it raises the average molecular speed only slightly.
        Explain. \workspace[2.8cm]
        \keyonly{\begin{apcanswer}The rate is governed by the
        \emph{fraction} of collisions carrying at least the activation
        energy, and that fraction sits far out in the tail of the
        distribution of molecular energies.

        Raising the temperature shifts and broadens the whole
        distribution. Near the average this is a small change, but out in
        the high-energy tail the population increases sharply --- a small
        shift can double or treble the number of molecules above the
        threshold. The increase in collision \emph{frequency} is minor by
        comparison; almost all the effect comes from the energy
        distribution.\end{apcanswer}}
\end{parts}

\begin{rubric}
  \rpoint{2}{(a) 1 pt each requirement}
  \rpoint{2}{(b) 1 pt the fraction above the activation energy rises
    sharply; 1 pt this outweighs the increase in collision frequency.
    ``Molecules move faster so they collide more'' alone earns 1}
\end{rubric}

% =====================================================================
\frq{short}{4}{Zumdahl \S12.1 \textbullet\ CED 5.1 \textbullet\ expressing reaction rate}

For \ce{2N2O5 -> 4NO2 + O2}, oxygen is produced at
\SI{0.0125}{\Molar\per\second}.

\begin{parts}
  \item Determine the rate at which \ce{NO2} is formed and the rate at
        which \ce{N2O5} is consumed. \workspace[2.6cm]
        \keyonly{\begin{apcanswer}Rates scale with the coefficients.

        \ce{NO2}: $4 \times 0.0125 = \mathbf{0.0500}$~M/s formed.

        \ce{N2O5}: $2 \times 0.0125 = \mathbf{0.0250}$~M/s consumed
        --- a \emph{decrease}, so
        $\Delta[\ce{N2O5}]/\Delta t = -0.0250$~M/s.\end{apcanswer}}
  \item State why a rate expressed as
        $-\tfrac{1}{2}\Delta[\ce{N2O5}]/\Delta t$ gives the same value as
        $\tfrac{1}{4}\Delta[\ce{NO2}]/\Delta t$.
        \answerlines[2]{Dividing each concentration change by its own
        coefficient gives a single rate for the reaction as a whole,
        independent of which species is monitored.}
\end{parts}

\begin{rubric}
  \rpoint{2}{(a) 1 pt \SI{0.0500}{\Molar\per\second} for \ce{NO2}; 1 pt
    \SI{0.0250}{\Molar\per\second} for \ce{N2O5} with the negative sign or
    the word ``consumed''}
  \rpoint{2}{(b) 1 pt dividing by coefficients; 1 pt gives one rate for
    the reaction regardless of species}
\end{rubric}

\examtotal

\end{document}
